% Problem dossier: VI.03.P6
\subsection*{Problem VI.03.P6 --- A tower of closed subsets and quotient maps}
\label{prob:vi03-p6}
\paragraph{Problem.}
Let \(Z\subseteq Y\subseteq X\subseteq\mathbb A_k^n\) be algebraic.  Describe the natural
surjections
\[
k[X]\twoheadrightarrow k[Y]\twoheadrightarrow k[Z]
\]
and prove that their composition is the natural quotient map \(k[X]\twoheadrightarrow k[Z]\).
Identify all three kernels in terms of vanishing ideals.
\ifdefined\FullProblemDossiers
\paragraph{Why this problem matters.} Contravariance must remain coherent under composition; this is the prototype for functoriality in VI/04.
\paragraph{Useful ideas before starting.} Use \(I(X)\subseteq I(Y)\subseteq I(Z)\).
\paragraph{Guided hints.} Write each ring as the same polynomial ring modulo a nested ideal.
\paragraph{Solution.}
The maps are induced by the identity on \(k[x_1,\dots,x_n]\):
\[
\frac{k[x]}{I(X)}\to\frac{k[x]}{I(Y)}\to\frac{k[x]}{I(Z)}.
\]
Their kernels are respectively \(I(Y)/I(X)\), \(I(Z)/I(Y)\), while the composite has
kernel \(I(Z)/I(X)\).  Since all maps send \([f]\) to the class of the same polynomial
\(f\), the composite is exactly the natural quotient map to \(k[Z]\).
\paragraph{What happened mathematically.} Nested geometric constraints become nested algebraic quotients, with arrows reversed.
\paragraph{Extensions.} Reinterpret the same diagram for regular maps after VI/04.
\paragraph{Provenance.} New synthesis problem emphasizing the arrow-reversal architecture.
\fi
